
\section{Expert review}

\section{The minimal viable product}
The minimal viable product (MVP) is a version of a new product that includes only the essential features necessary to meet the needs for initial testing. The goal of an MVP is to validate the product idea with minimal resources and gather feedback for future development. In figure \ref{fig:MVP} you can see the result of the MVP and in figure \ref{fig:MVPpipeline} you can see the pipeline that is used for the MVP.

\begin{figure}[!htp]
    \centering
    \includegraphics[scale=0.1]{images/mvpTTRPGtoComic.png}
    \caption[The minimal viable product]{The minimal viable product}
    \label{fig:MVP}
\end{figure}

\subsection{The minimal viable product}

\section{Small experiments}

\subsection{The language of TTRPGs}
Models like Whisper are trained on a wide variety of audio data, however this is primarily normal language usage. Because of this I anticipate that there will be a large issue with transcribing many fantasy words and terms, thus a small experiment has been done to identify if this is truly a problem. For this experiment 10 fantasy creature names have been recorded and whisper transcribed these 10 names. Only 1 out of the 10 names were transcribed correctly and 2 out of 10 where almost correct \ref{tab:monster_names}. This is inaccuracy is a big issue when transcribing any fantasy game session as it will interfere with many important key moments and key names during a session/campaign. When the group meets a monster called a "Slaad" and Whisper transcribes this as "Nod" it will quite confusing to understand what is going on in the story.

To generate a accurate transcription of the names a custom model needs to be trained on specific fantasy words and terms, however this is way outside of the scope of this thesis. In addition it is also very difficult as Fantasy names are fantasy names for a reason and people can just come up with new names all the time. Integrating thousands of names however could still be valuable as it will most likely catch most cases.
\begin{table}[h!]
\centering
\begin{tabular}{|l|l|}
\hline
\textbf{Correct Name} & \textbf{Transcribed name} \\
\hline
Owlbear    & Halber     \\
Goblin     & Goblin     \\
Rakshasa   & Raksasa    \\
Dracolich  & Draco Lich \\
Quasit     & Kwasi      \\
Slaad      & Nod        \\
Drow       & Drought    \\
Illithid   & Illidan    \\
Cambion    & Kempion    \\
Ankheg     & Hank Man   \\
\hline
\end{tabular}
\caption{Comparison of monster names and whisper transcribed versions}
\label{tab:monster_names}
\end{table}


%This is old, might be deleted
\section{Generating the comic book script}
The iterative design process of generating the comic book script after the MVP version was a crucial step to improve the quality of the comic book panels. From having a one-shot approach where we gave the entire transcript to the LLama 3.1 model, we moved to an adaptive panel generator which has a structured approach to data and keeps track of the context. In this section all the iterative steps are described and what did and did not work in chronological order.

\subsubsection{MVP}
For the MVP of the artefact the entire transcript of the tabletop role-playing game session was given to Llama 3.1 and see what it would generate. The LLM was asked: "You are given a Dungeons and Dragons session transcript. Convert it into a comic book script." This resulted in two results: The LLM crashed as the token limit was exceeded or it gave things that had nothing to do to with the input and hallucinated everything. A quick fix to this was to cut the transcript into smaller chunks of text which resulted in an interesting result as it actually gave a comic book script. It was given 2 scenes of text which where hand cut out of the transcript. However as the LLM was not instructed to give a structured output it was useless as every time the program was run it would give a different output. The last step was to create a structured prompt for the LLM which is the prompt in \ref{lst:MPVLlamaprompt}. This resulted in a consistent JSON format output which could then be used to generate images.

\subsubsection{Cleaning up the transcript}
When you feed an audio recording to Whisper it transcribes the audio as one long sentence. There is no punctation or any idea of sentences structure. This way it is impossible to keep track of who said what and when. To fix this the transcript needs to be cleaned up. The first step is to use a Punctuation and Capitalization Restoration model \cite{guhrEtAl} to add punctuation and capitalization to the text. This will make the text more readable and easier to understand for an LLM. The second step is to use Speaker Diarization to identify the different speakers in the text. The tool used to do this was pyannote 3.1 \cite{Bredin23,Plaquet23}, an open source python library which is super easy to use and broadly used by people (15 million download per month). Lastly combine the cleaned up whisper transcription with the speaker diarization to create a structured JSON file. This new transcript file \ref{lst:sessionTrarnscriptJSON} is what created the foundation for the entire comic book generator.

\begin{lstlisting}[style=jsonstyle, caption={Session transcript JSON}, label={lst:sessionTrarnscriptJSON}]
{
    "panel_index": 0,
    "speaker": "SPEAKER_02",
    "text": "Our story actually starts with the three of you...",
    "start": 0.0,
    "end": 2.96
},
...
\end{lstlisting}

\subsubsection{Scene detection}

\subsubsection{Line based chunking}
Every line or every X lines

\subsubsection{LLM based chunking}
Scene cutting

\subsubsection{LLM based filtering}
Ask the LLM to decide if a line of text is player or character speech. Into giving it a chunk instead.

\subsubsection{Context aware LLM based chunking}

\subsubsection{Adaptive panel generating}